\subsection{ARCH-136-SKELETON-OBSTRUCTION: a common finite skeleton does not give arbitrary relative completion}

\paragraph{Archive and scope.}
This tests the higher-simplex step after the accepted common first
skeleton, for cote 136, physical p.~121, footer 120,
\texttt{ARCHIVE\_136\_scan121.tex}, with the uniform tests on
physical pp.~117 and 119, footers 116 and 118, scans 117 and 119.
It is a shifted form of the supplied AQ-136-HIGHER-COMPLETION
argument, proved here because the preservation of a specified
finite skeleton is an additional assertion.  It is not a
counterexample to the original common-source question.

\paragraph{Statement.}
For every $d\geq1$, prime $p$, and $m\geq d+1$, there is a
simplicial strict differential PD $S$-algebra $A$ and two maps
$f_0,f_1:A\rightrightarrows B$ to the standard admissible theory,
with
\[
 \begin{gathered}
 A_n=B_n\quad(n<d),\qquad
 f_0|_{\operatorname{sk}_{d-1}A}
     =f_1|_{\operatorname{sk}_{d-1}A}=1,\\
 H^0(A_n^{w,*})=S^w,\qquad
 H^j(A_n^{w,*})=0\quad(1\leq j\leq d),\\
 H^*((A_n)_{\mathbf Q})=S_{\mathbf Q},\qquad
 H_r(NA_*^{w,q})=0\quad(w>mq,\ r\geq0),
 \end{gathered}                                           \tag{136.SO1}
\]
but no factorization
\[
       A\longrightarrow R
          \mathrel{\substack{\xrightarrow{j_0}\\[-3pt]
                            \xrightarrow[j_1]{}}}B,
       \qquad R\in\mathcal A_{\mathbf Z},\quad
                    j_i\circ(A\to R)=f_i               \tag{136.SO2}
\]
exists.  In particular $d=2$ gives exact agreement on the entire
first skeleton and exact vertical cohomology through degree two,
yet still obstructs the prescribed simultaneous completion in
arbitrarily high weights.

\paragraph{Proof.}
Let $U$ be the Dold--Kan image of $\mathbf Z[d]$.  Make copies
$V,W$ in form degrees $d,d+1$, both of weight $m$, and take the
free graded PD algebra
\[
 D=\begin{cases}
       \Gamma(V)\otimes\Lambda(W),&d\text{ even},\\
       \Lambda(V)\otimes\Gamma(W),&d\text{ odd},
    \end{cases}
 \qquad dv=p u,\quad du=0,
 \qquad A=B\otimes_{\mathbf Z}D.                         \tag{136.SO3}
\]
Here $v\mapsto u$ is the identity of the two copies of $U$, so
the construction is simplicial.  The free graded PD derivation
has square zero and preserves weights.  Because $U_n=0$ for $n<d$,
$A_n=B_n$ in exactly the asserted skeleton.

In each simplex $U_n$ is free, so $D_n$ is a tensor product of
free disks with differential multiplication by $p$.  Its first
positive term is in degree $d$, and the differential from that
term to the independent degree-$(d+1)$ generators is injective.
Therefore $H^j(D_n)=0$ for $1\leq j\leq d$ and $H^0(D_n)=\mathbf Z$.
The integral contraction of $B_n$ onto $S$ shows the same assertions
for $A_n$.  Over $\mathbf Q$ each disk contracts, including the
odd case where the other block multipliers are positive integers
times $p$.  This proves the rational assertion.

A monomial of $D$ of form degree $q'$ and weight $e$ satisfies
$e\leq m q'$.  Thus, if $w>mq$, every summand
$B^{w-e,q-q'}\otimes D^{e,q'}$ has
$w-e>m(q-q')\geq q-q'$.  Its $B$ factor is simplicially
contractible.  Tensor its linear homotopy with the identity on
the other factor.  There are only finitely many possible $(e,q')$
for a fixed $q$, and their sum proves the horizontal assertion
of (136.SO1).

Let $x_0,\ldots,x_d$ and $e_i=dx_i$ be the standard coordinates
on $B_d$.  Define the top form
\[
                 \eta=x_0^{[m-d]}e_1\cdots e_d
                          \in B_d^{m,d}.                 \tag{136.SO4}
\]
Its restriction to every face is zero because a $(d-1)$-simplex
has no form of degree $d$; it is vertically closed because $B_d$
has no form of degree $d+1$.  It therefore gives a simplicial
linear map $U\to B^{m,d}$.  On $D$ define $f_0(v)=f_0(u)=0$,
$f_1(v)=\eta$ and $f_1(u)=0$, and make both maps the identity
on $B$.  The free differential PD property proves that these
are strict maps.  They agree on the displayed finite skeleton.
Notice that $m>d$ even puts this chosen input in the acyclic
form-$d$ row of $B$; the nonunit relation $dv=pu$ remains material.

Put $X=\Delta[d]/\partial\Delta[d]$.  The normalized generators
$v,u$ have every face zero, so define cochains on this sphere.
For $w\geq m$ write
\[
       v_w=T^{[w-m]}v,\quad u_w=T^{[w-m]}u,
                              \qquad dv_w=p u_w.          \tag{136.SO5}
\]
The class of $f_1(v_w)$ in the comparison range is
\[
            \binom wm T^{[w]}
                         \in H^d(B^{w,*}(X))\simeq S^w.  \tag{136.SO6}
\]
Here is the normalization check.  Integral PD simplex integration
has the basis formula
\[
 \int_{\Delta[d]}x_0^{[r_0]}\cdots x_d^{[r_d]}
                      e_1\cdots e_d
                           =T^{[r_0+\cdots+r_d+d]}.
\]
It follows from the ordinary polynomial simplex integral over
$\mathbf Q$ by cancellation of the factorials; both sides are
integral divided-power basis elements.  Stokes' identity holds
over $\mathbf Q$, hence integrally, so this defines a cochain map
to ordinary integral simplicial cochains with coefficients $S^w$.
On the sphere, the form $x_0^{[w-d]}e_1\cdots e_d$ has integral
one, so in the range where AQ-136-COMPARISON gives $H^d\simeq\mathbf Z$
this integral is an isomorphism onto $S^w$.  Formula (136.SO4)
has integral $T^{[m]}$; multiplying by $T^{[w-m]}$ gives exactly
(136.SO6), by the coefficient PD multiplication rule.

Suppose (136.SO2) existed, and denote $A\to R$ by $g$.
Its restriction $h:B\to R$ satisfies $j_0h=j_1h=1_B$.
In sufficiently large weights the cohomological comparison for
$B,R$ on $X$ makes $H^d(h)$ an isomorphism.  Hence $H^d(j_0)$
and $H^d(j_1)$ are equal, both its inverse.  The same comparison
gives $H^{d+1}(R^{w,*}(X))=0$.  Choose $a$ of degree $d$ with
$da=g(u_w)$.  Then $pa-g(v_w)$ is closed.  Its two images have
equal degree-$d$ cohomology classes, while each $j_i(a)$ is closed
because $f_i(u_w)=0$.  Thus (136.SO6) forces
\[
 p\bigl([j_1(a)]-[j_0(a)]\bigr)
                              =\binom wm T^{[w]}
                                    \quad\text{in }S^w.  \tag{136.SO7}
\]
There are arbitrarily large $w$ for which the right side is not
divisible by $p$.  Specifically take $w=m+p^r$ with $p^r>m$.
The identity $(1+Z)^{p^r}=1+Z^{p^r}$ in $\mathbf F_p[Z]$ shows
that $\binom{m+p^r}{m}=1$ modulo $p$.  Such a weight can exceed
every fixed comparison bound in the argument.  Equation (136.SO7)
is then impossible.  This proves the obstruction.\hfill$\square$

\paragraph{Implication for the common-first-skeleton route.}
Agreement on a common first skeleton makes the degree-two vertex
residues identical, because they can be read on the common edges.
It does not ensure that all subsequent degree-three repairs can
preserve that agreement.  The $d=2$ example has exactly that common
first skeleton and nevertheless has the next residue on a
simplicial two-sphere.  The same construction moves the obstruction
above any prescribed finite skeleton.  It excludes an unrestricted
relative completion theorem based only on finite-skeleton agreement;
it leaves open whether the specific free construction can maintain
additional compatible data throughout its repairs.
